\section*{Bab \ref{ch:b3:probability} --- Peluang:
Landasan dan Hukum Bilangan Besar}

\begin{solution}{exo:b3:probability:1}
(a) Untuk $u \in \intoo01$: $\P(F(X) \leq u) = \P(X \leq
F^{-1}(u)) = F(F^{-1}(u)) = u$ (sebab \omterm{def:b3:topology:continuity}{kekontinuan} dan kemonotonan
sejatinya menjadikan $F$ bijeksi ke $\intoo01$ dengan
$\{F(X) \leq u\} = \{X \leq F^{-1}(u)\}$): jadi $F(X)$ seragam.
Sebaliknya $\P(G(U) \leq t) = \P(U \leq F(t)) = F(t)$: jadi untuk
menyimulasikan sebuah \omterm{def:b3:probability:space}{distribusi}, terapkanlah fungsi \omterm{def:b3:probability:space}{distribusi} balikannya pada
cuplikan seragam.

(b) Ambillah $Y = X^2$ dengan $X$ seragam pada $\intcc{-1}1$: bagi $t \in
\intcc01$, $F_Y(t) = \P(-\sqrt t \leq X \leq \sqrt t) = \sqrt
t$: jadi \omterm{ex:b3:lebesgue:gamma}{kepadatannya} $\frac1{2\sqrt t}\mathbf 1_{\intoo01}$. Sedangkan
$\P\bigl(-\frac1\lambda\ln U \leq t\bigr) = \P(U \geq
\eu^{-\lambda t}) = 1 - \eu^{-\lambda t}$: yakni eksponensial
$\mathcal E(\lambda)$ --- yakni pembalikan dalam aksi.
\end{solution}

\begin{solution}{exo:b3:probability:2}
(a) Poisson: $\E X = \sum_{k\geq0}k\,\eu^{-\lambda}
\frac{\lambda^k}{k!} = \lambda$, $\E[X(X-1)] = \lambda^2$, sehingga
$\V = \lambda^2 + \lambda - \lambda^2 = \lambda$. Geometrik
(dengan $\P(X = k) = p(1-p)^{k-1}$): $\E X = \frac1p$, $\V =
\frac{1-p}{p^2}$ (turunkanlah deret geometriknya dua kali).

(b) Fungsi $G(t) = \P(T > t)$ tak naik dengan $G(0^+)\dots$
$G \colon \intco0\infty \to \intoc01$; dan ketiadaan ingatannya berbunyi
$G(t + s) = G(t)G(s)$. Maka $G(n t) = G(t)^n$ dan $G(t/n) =
G(t)^{1/n}$: sehingga $G(q) = G(1)^q$ bagi $q \geq 0$ yang rasional; lalu dengan menulis
$G(1) = \eu^{-\lambda}$ (dengan $\in \intoo01$: sebab $G(1) = 1$ akan
memaksa $G \equiv 1$, yang mustahil bagi \omterm{def:b3:probability:space}{peubah acak} berhingga;
sedangkan $G(1) = 0$ terkecualikan oleh hipotesisnya) lalu mengapit
sembarang $t$ di antara bilangan rasional (menurut kemonotonannya): $G(t) =
\eu^{-\lambda t}$ --- yakni \omterm{def:b3:probability:space}{distribusi} eksponensialnya. Sedangkan konversnya sebuah
perhitungan.
\end{solution}

\begin{solution}{exo:b3:probability:3}
(a) Berlaku $\sigma(f_i(X_i)) = f_i(X_i)^{-1}(\mathcal B) \subseteq
X_i^{-1}(\mathcal B) = \sigma(X_i)$ (sebab $f_i$ Borel), sedangkan
sub-$\sigma$-aljabar dari $\sigma$-aljabar yang \omterm{def:b3:probability:independence}{bebas} juga
\omterm{def:b3:probability:independence}{bebas} (sebab identitas pendefinisinya berlaku lebih-lebih lagi).

(b) Berlaku $\sigma(A_i) = \{\varnothing, A_i, A_i^c, \Omega\} =
\sigma(A_i^c) = \sigma(\mathbf 1_{A_i})$: sehingga ketiga
pernyataannya menegaskan \omterm{def:b3:probability:independence}{kebebasan} $\sigma$-aljabar yang
sama. (Bahwa pemfaktoran atas $A_i$
merambat ke komplemennya adalah argumen $\lambda$-sistem
di dalam kesetaraan \cref{def:b3:probability:independence}
--- atau inklusi-eksklusi langsung.)

(c) Berlaku $\P(A) = \P(B) = \P(C) = \frac12$; sedangkan $A\cap B = A\cap C =
B\cap C$ pada pasangannya: sebab setiap irisannya adalah ``keduanya gambar'' atau
yang serupa, berpeluang $\frac14$: jadi \omterm{def:b3:probability:independence}{bebas} berpasangan.
Padahal $\P(A\cap B\cap C) = \P(\text{GG}) = \frac14 \neq
\frac18$: jadi tidak \omterm{def:b3:probability:independence}{bebas} --- sebab $C$ ditentukan oleh $A$ dan
$B$.
\end{solution}

\begin{solution}{exo:b3:probability:4}
(a) Misalkan teks $T$ berpanjang $L$ dan $q = a^{-L}$ (dengan $a$
sebagai \omterm{def:b3:measure:measure}{ukuran} abjadnya). Kejadian $E_k = \{$posisi $kL+1,
\dots, (k+1)L$ mengeja $T\}$ saling \omterm{def:b3:probability:independence}{bebas} (sebab blok huruf i.i.d.\ yang
saling lepas), masing-masing berpeluang $q > 0$: jadi $\sum\P(E_k)
= \infty$, sehingga Borel--Cantelli (2) memberikan tak berhingga banyak
kemunculan secara h.p.

(b) Untuk yang atas: $\P\bigl(R_n \geq (1+\varepsilon)\log_2n\bigr)
\leq 2^{-(1+\varepsilon)\log_2n} = n^{-(1+\varepsilon)}$,
yang terjumlahkan: jadi menurut Borel--Cantelli (1), secara h.p.\ hanya berhingga banyak
$n$ semacam itu. Untuk yang bawah: kemaslah blok saling lepas --- yang ke-$j$
berpanjang $\ell_j = \lceil\log_2s_j\rceil$ dan mulai di $s_j =
\sum_{i<j}\ell_i$; lalu kejadian ``blok $j$ seluruhnya satu'' saling
\omterm{def:b3:probability:independence}{bebas} dengan peluang $2^{-\ell_j} \asymp
\frac1{s_j} \asymp \frac1{j\log_2 j}$, yang jumlahnya divergen:
sehingga Borel--Cantelli (2) memberikan tak berhingga banyak blok seluruhnya satu,
yakni $R_{s_j} \geq \log_2 s_j$ tak berhingga kali. Bersama-sama:
panjang rentetan maksimal pada $n$ digit pertamanya adalah
$(1 + o(1))\log_2n$ secara h.p.
\end{solution}

\begin{solution}{exo:b3:probability:5}
(a) Nilai $R = \bigl(\limsup\abs{X_n}^{1/n}\bigr)^{-1}$ tak
berubah jika berhingga banyak $X_n$ diubah: sehingga bagi setiap $N$,
$R$ terukur-$\sigma(X_N, X_{N+1}, \dots)$, yakni
terukur-ekor. Maka setiap kejadian $\{R \leq c\}$ berpeluang
$0$ atau $1$ (\cref{thm:b3:probability:zeroone}),
sehingga fungsi \omterm{def:b3:probability:space}{distribusi} $R$ hanya mengambil nilai
$0, 1$: jadi ia melompat pada satu titik $c_0 \in
\intcc0{+\infty}$, dan $R = c_0$ secara h.p.

(b) Konvergensi $\sum X_n$ dan konvergensi $\frac{S_n}n$ tak
peka terhadap perubahan berhingga banyak sukunya (untuk yang kedua:
sebab suku yang diubah menyumbang $O(1/n) \to 0$): jadi keduanya kejadian ekor;
lalu \omterm{thm:b3:probability:zeroone}{hukum nol--satunya} berlaku.

(c) Kejadian $\{X_1 > 0\}$ bergantung pada $X_1$: bagi tanda i.i.d.\
(dengan $\P(X_1 = \pm1) = \frac12$), peluangnya $\frac12
\notin \{0,1\}$ --- dan tak ada pertentangan, sebab ia bukan kejadian
ekor.
\end{solution}

\begin{solution}{exo:b3:probability:6}
Bekerjalah pada $(\intcc01, \lambda)$.
(a) Mesin tiknya $\mathbf 1_{I_n}$
(\cref{exo:b3:lp:3}): $\norm{X_n}_p^p = \lambda(I_n) \to 0$
(bagi setiap $p < \infty$), sehingga juga dalam peluang; padahal pada setiap
$\omega$ nilai $0$ dan $1$ keduanya berulang: jadi tak ada konvergensi
titik demi titik di mana pun.
(b) Berlaku $X_n = n\mathbf 1_{\intoo0{1/n}} \to 0$ di luar $0$, tetapi
$\norm{X_n}_p \geq n^{1 - 1/p} \geq 1$.
(c) Ambillah $X_n = \sqrt n\,\mathbf 1_{\intoo0{1/n}}$: maka $\E\abs{X_n} =
n^{-1/2} \to 0$, sedangkan $\E X_n^2 = 1$.
(d) Dari sembarang subbarisan sarikanlah (menurut konvergensi dalam
peluangnya) subbarisan lebih lanjut yang konvergen secara h.p.\
(\cref{prop:b3:probability:modes}(c)); lalu konvergensi terdominasinya
memberikan konvergensi $L^1$ sepanjangnya, dengan limit yang \emph{sama}
$X$. Jadi setiap subbarisan barisan numerik
$\E\abs{X_n - X}$ mempunyai subsubbarisan yang menuju $0$: sehingga
seluruh barisannya menuju $0$.
\end{solution}

\begin{solution}{exo:b3:probability:7}
(a) Berlaku $\hat p_n = \frac{S_n}n$ dengan $S_n$ yang binomial: sehingga $\V(\hat
p_n) = \frac{p(1-p)}n \leq \frac1{4n}$, lalu Chebyshev
(\cref{prop:b3:probability:markov}) memberikan batasnya.
(b) Selesaikanlah $\frac1{4n(0.03)^2} \leq 0.05$: $n \geq
\frac{1}{4\cdot0.0009\cdot0.05} \approx 5556$. Sedangkan teorema
limit pusatnya akan membenarkan $n \approx 1070$ bagi jaminan yang sama:
jadi Chebyshev membayar keumumannya dengan faktor
$\approx 5$.
\end{solution}

\begin{solution}{exo:b3:probability:8}
(a) Jika $S_n \sim \mathcal B(n, x)$, maka teorema transfernya
memberikan $\E\bigl[f(\frac{S_n}n)\bigr] =
\sum_k\binom nkx^k(1-x)^{n-k}f(\frac kn) = B_nf(x)$.

(b)--(c) Misalkan $\omega = \omega_f$ modulus \omterm{def:b3:topology:continuity}{kekontinuannya}
(dengan $\abs{f(u) - f(v)} \leq \omega(\abs{u - v})$, dan $\omega(c
\delta) \leq (1 + c)\,\omega(\delta)$ lewat perantaian langkahnya).
Maka, bagi sembarang $\delta > 0$,
\[
\abs{f(u) - f(x)} \leq \Bigl(1 + \frac{(u -
x)^2}{\delta^2}\Bigr)\omega(\delta)
\]
(jika $\abs{u - x} \leq \delta$, jelas; sedangkan jika tidak
$\omega(\abs{u-x}) \leq (1 + \frac{\abs{u-x}}\delta)
\omega(\delta) \leq (1 + \frac{(u-x)^2}{\delta^2})
\omega(\delta)$). Lalu ambillah \omterm{def:b3:probability:space}{nilai harapannya} di $u = \frac{S_n}n$:
\[
\abs{B_nf(x) - f(x)} \leq
\Bigl(1 + \frac{\V(S_n/n)}{\delta^2}\Bigr)\omega(\delta)
\leq \Bigl(1 + \frac{1}{4n\delta^2}\Bigr)\omega(\delta) ;
\]
dengan $\delta = n^{-1/2}$: $\norm{B_nf - f}_\infty \leq
\frac54\,\omega\bigl(n^{-1/2}\bigr) \leq
\frac32\,\omega\bigl(n^{-1/2}\bigr) \to 0$ (menurut \omterm{def:b3:topology:continuity}{kekontinuan}
seragamnya pada \omterm{def:b3:topology:compact}{kompak}): yakni teorema Weierstrass yang bersifat peluang,
dengan laju yang eksplisit dan seragam.
\end{solution}

\begin{solution}{exo:b3:probability:9}
(a) Setelah $k - 1$ jenis terkumpul, setiap penarikan bersifat baru dengan
peluang $p_k = \frac{n-k+1}n$: jadi $\tau_k$ geometrik
$(p_k)$, dan $\tau_k$ saling \omterm{def:b3:probability:independence}{bebas} (sebab penarikannya \omterm{def:b3:probability:independence}{bebas}).
Lalu jumlahnya: $\E T_n = \sum_k\frac n{n-k+1} = nH_n \sim n\ln n$;
$\V(T_n) = \sum\frac{1 - p_k}{p_k^2} \leq
n^2\sum_{j=1}^n\frac1{j^2} \leq \frac{\pi^2}6n^2$.

(b) Chebyshev: $\P\bigl(\abs{T_n - nH_n} \geq \varepsilon
n\ln n\bigr) \leq \frac{\pi^2n^2/6}{\varepsilon^2n^2\ln^2n}
\to 0$, sedangkan $\frac{nH_n}{n\ln n} \to 1$: sehingga $\frac{T_n}{n\ln n}
\to 1$ dalam peluang.

(c) Lewat batas gabungannya: $T_n > t$ berarti suatu jenis tak terlihat setelah
$\lceil t\rceil$ penarikan, sehingga $\P(T_n > t) \leq n(1 -
\frac1n)^{t} \leq n\,\eu^{-t/n}$; lalu di $t = \beta n\ln n$:
$\leq n^{1 - \beta}$. Untuk $\beta > 1$, $\sum_m
2^{m(1-\beta)} < \infty$: sehingga Borel--Cantelli memberikan, sepanjang $n =
2^m$, secara h.p.\ $T_n \leq \beta n\ln n$ pada akhirnya --- khususnya
bagi setiap $\beta > 2$ seperti yang dinyatakan (dan sembarang $\beta > 1$
berhasil sepanjang subbarisannya).
\end{solution}

\begin{solution}{exo:b3:probability:10}
(a) Digit berindeks genap $(b_{2k})_k$ merupakan bit adil i.i.d.\
(yakni subkeluarga keluarga digitnya yang \omterm{def:b3:probability:independence}{bebas}), sehingga $U =
\sum_kb_{2k}2^{-k}$ memberi setiap selang diadik peluangnya yang benar
(seperti pada
\cref{thm:b3:probability:existence}): jadi seragam; demikian pula $V$;
sedangkan $(U, V)$ bergantung pada blok digit yang saling lepas: jadi \omterm{def:b3:probability:independence}{bebas}
(lewat pemfaktoran pada persegi panjang diadiknya, lalu Dynkin).

(b) Pemetaan $\Phi(\omega) = (U(\omega), V(\omega))$ \omterm{def:b3:lebesgue:measurable}{terukur}
dengan $\Phi_*\lambda = \lambda\otimes\lambda = \lambda_2$
(menurut kesamaannya pada persegi panjang diadik + ketunggalannya). Lalu menyelang-nyelingkan
digitnya mendefinisikan sebuah balikan yang terdefinisi di luar himpunan (nol) berisi
bilangan rasional diadik pada salah satu faktornya: yakni bijeksi
yang mengawetkan \omterm{def:b3:measure:measure}{ukuran} antara himpunan bagian berukuran penuh $\intcc01$ dan
$\intcc01^2$. Bandingkanlah dengan Peano
(\cref{pb:b3:topology:1}): sebab \omterm{def:b3:topology:continuity}{kekontinuannya} memaksa kesurjektifan
tanpa keinjektifan; sedangkan menukar \omterm{def:b3:topology:continuity}{kekontinuan} dengan sekadar
\omterm{def:b3:lebesgue:measurable}{keterukuran} membeli \omterm{def:b3:measure:measure}{isomorfisma-ukuran} --- jadi dimensi
tak terlihat oleh teori \omterm{def:b3:measure:measure}{ukuran}, tetapi terlihat oleh \omterm{def:b3:topology:topology}{topologi}.
\end{solution}

\begin{solution}{exo:b3:probability:11}
(a) \omterm{def:b3:topology:continuity}{Kekontinuan} \omterm{def:b3:probability:space}{distribusinya} menjadikan seri sebagai kejadian nol
(seperti pada argumen statistik terurut di babnya), sedangkan
$n!$ pengurutan relatif $(X_1, \dots, X_n)$ bersifat
tertukarkan, sehingga sama mungkinnya. Lalu $R_n = 1$ berarti
maksimumnya duduk pada posisi terakhir: jadi berpeluang
$\frac{(n-1)!}{n!} = \frac1n$.

(b) Tetapkanlah $n$ lalu syaratkan pada urutan relatif $X_1,
\dots, X_{n-1}$: menyisipkan $X_n$ ke $n$ celah peringkat
yang mungkin bersifat seragam dan \omterm{def:b3:probability:independence}{bebas} dari urutan itu
(menurut ketertukaran $n$-tupelnya). Karena itu $R_n$ (yakni kejadian
``$X_n$ menempati celah teratas'') \omterm{def:b3:probability:independence}{bebas} dari seluruh
riwayat rekornya $(R_1, \dots, R_{n-1})$, yang merupakan fungsi
urutan relatif $n - 1$ variabel pertamanya.
Lalu induksinya memberikan \omterm{def:b3:probability:independence}{kebebasan} penuh dengan $\P(R_n = 1) =
\frac1n$.

(c) Berlaku $\sum\P(R_n = 1) = \sum\frac1n = \infty$ beserta
\omterm{def:b3:probability:independence}{kebebasannya}: sehingga paruh kedua Borel--Cantelli memberikan rekor
tak berhingga kali secara h.p.\ (rekornya tak pernah berhenti --- tetapi
menipis secara logaritmik: $\E[\#\text{rekor} \leq n] =
H_n \approx \ln n$). Untuk rekor berurutan: $\P(R_n = R_{n+1}
= 1) = \frac1{n(n+1)}$ (menurut \omterm{def:b3:probability:independence}{kebebasannya}), dan
\[
\sum_n\frac1{n(n+1)} = \sum_n\Bigl(\frac1n -
\frac1{n+1}\Bigr) = 1 < \infty :
\]
sehingga paruh pertama Borel--Cantelli berlaku --- jadi hanya berhingga
banyak pasangan rekor berurutan yang terjadi, secara h.p.
\end{solution}

\begin{solution}{exo:b3:probability:12}
(a) Sebuah rentetan berpanjang $\ell$ yang mulai di posisi $i \leq n$
berpeluang $2^{-\ell}$; lalu lewat batas gabungannya: $\P(L_n \geq \ell)
\leq n2^{-\ell}$. Dengan $\ell_n = (1 +
\varepsilon)\log_2n$: $\P(L_n \geq \ell_n) \leq
n^{-\varepsilon}$. Lalu sepanjang $n = 2^k$:
$\sum_k2^{-k\varepsilon} < \infty$, sehingga secara h.p.\ $L_{2^k} <
(1+\varepsilon)k$ pada akhirnya (menurut Borel--Cantelli); sedangkan bagi $n$
yang umum ambillah $2^{k-1} < n \leq 2^k$ lalu pakailah kemonotonan
$L_n$ ditambah $\log_22^{k-1} \leq \log_2n$: $L_n \leq L_{2^k}
< (1 + \varepsilon)k \leq (1 + \varepsilon)\frac{k}{k-1}
\log_2n$, dan faktor tambahannya terserap dengan sedikit membesarkan
$\varepsilon$.

(b) Dengan $\ell = \lceil(1 - \varepsilon)\log_2n\rceil$ dan
$m = \lfloor n/\ell\rfloor$ blok saling lepas: bloknya saling
\omterm{def:b3:probability:independence}{bebas}, masing-masing seluruhnya gambar dengan peluang $2^{-\ell}
\geq n^{-(1-\varepsilon)}/2$, sehingga
\[
\P(L_n < \ell) \leq \bigl(1 - 2^{-\ell}\bigr)^{m}
\leq \exp\bigl(-m2^{-\ell}\bigr)
\leq \exp\Bigl(-c\,\frac{n^{\varepsilon}}{\log_2n}\Bigr)
\]
bagi sebuah konstanta $c > 0$ dan $n$ yang besar. Peluang ini
terjumlahkan sepanjang $n = 2^k$ (bahkan sepanjang semua $n$):
sehingga Borel--Cantelli memberikan secara h.p.\ $L_n \geq (1 -
\varepsilon)\log_2n$ pada akhirnya (sedangkan kemonotonannya mengisi
di antara $2^k$ seperti pada (a), tanpa bahaya).

(c) Kedua batasnya sepanjang barisan $\varepsilon = \frac1j$, lalu
dengan mengiriskan terbilang banyak kejadian berukuran penuh:
$\frac{L_n}{\log_2n} \to 1$ secara h.p. Untuk $n = 10^6$: $\log_2n
\approx 19.9$ --- jadi rentetan $\approx 20$ gambar bukanlah
anomali yang mencurigakan melainkan kepastian matematis, sedangkan
ketiadaannya justru bukti bahwa manusia memalsukan ``keacakan'' (sebab manusia
jarang berani menulis lebih dari $5$ atau $6$ gambar berturut-turut).
\end{solution}

\begin{solution}{pb:b3:probability:1}
\textbf{1.} Variabel $X_n^{\pm}$ merupakan fungsi Borel $X_n$: sehingga keduanya
tetap \omterm{def:b3:probability:independence}{bebas} berpasangan
(\cref{exo:b3:probability:3}(a)) dan berdistribusi
identik, \omterm{def:b3:lebesgue:l1}{terintegralkan}, dengan $\E X_1 = \E X_1^+ - \E
X_1^-$. Lalu jika teoremanya berlaku bagi variabel tak negatif,
terapkanlah pada kedua paruhnya lalu kurangkan:
$\frac{S_n}n = \frac{S_n^+}n - \frac{S_n^-}n \to \E X_1^+ -
\E X_1^- = m$ secara h.p.

\textbf{2.} Berlaku $\P(X_n \neq Y_n) = \P(X_n > n) = \P(X_1 > n)$
(menurut \omterm{def:b3:probability:space}{distribusinya} yang identik), dan $\sum_n\P(X_1 > n) \leq \sum_n\P(X_1
\geq n) \leq \E X_1 < \infty$
(\cref{exo:b3:product:3}(a)). Lalu Borel--Cantelli (1): secara h.p.\
$X_n = Y_n$ bagi setiap $n$ yang besar, sehingga $S_n - S_n^*$
akhirnya konstan terhadap $n$: jadi $\frac{S_n - S_n^*}n \to 0$
secara h.p., dan kedua jumlah ternormalkannya berbagi perilaku
asimtotiknya.

\textbf{3.} Berlaku $X_1\mathbf 1_{X_1 \leq n} \nearrow X_1$: sehingga teorema
konvergensi monotonnya memberikan $\E Y_n \to m$; sedangkan rata-rata Cesàro
sebuah barisan konvergen menuju limit yang sama: $\frac{\E S_n^*}n =
\frac1n\sum_{k\leq n}\E Y_k \to m$. Karena itu cukup
dibuktikan $\frac{S^*_n - \E S^*_n}{n} \to 0$ secara h.p.

\textbf{4.} Berlaku $\V(Y_n) \leq \E Y_n^2 = \E[X_1^2\mathbf
1_{X_1\leq n}]$. Lalu menurut Tonelli bagi deret,
\[
\sum_n\frac{\E[X_1^2\mathbf 1_{X_1\leq n}]}{n^2}
= \E\Bigl[X_1^2\!\!\sum_{n \geq \max(X_1, 1)}\!\frac1{n^2}
\Bigr]
\leq \E\Bigl[X_1^2\cdot\frac{4}{\max(X_1,1)}\Bigr]
\leq 4\,\E[X_1] < \infty,
\]
dengan memakai $\sum_{n\geq x}n^{-2} \leq \frac4x$ bagi $x \geq 1$
(sebab bagi $x \geq 2$: $\leq \frac1{x-1} \leq \frac2x$; sedangkan bagi $1
\leq x < 2$: $\leq \frac{\pi^2}6 \leq \frac4x$ karena
$\frac4x > 2$), dan $X_1^2/\max(X_1, 1) \leq X_1$ pada kedua
kasus $X_1 \gtrless 1$.

\textbf{5.} \omterm{def:b3:probability:independence}{Kebebasan} berpasangannya memberikan $\E[(Y_i - \E
Y_i)(Y_j - \E Y_j)] = 0$ bagi $i \neq j$ (yakni rumus hasil kali
bagi dua variabel), sehingga variansnya menjumlah: $\V(S^*_k) =
\sum_{n\leq k}\V(Y_n)$. Lalu Chebyshev pada setiap $k_j$ lalu
menjumlahkannya:
\[
\sum_j\P\Bigl(\abs{S^*_{k_j} - \E S^*_{k_j}} \geq
\varepsilon k_j\Bigr)
\leq \frac1{\varepsilon^2}\sum_j\frac1{k_j^2}\sum_{n\leq
k_j}\V(Y_n)
= \frac1{\varepsilon^2}\sum_n\V(Y_n)\!\!\sum_{j : k_j\geq
n}\!\frac1{k_j^2}
\]
(menurut Tonelli bagi deret ganda tak negatifnya).

\textbf{6.} Berlaku $k_j = \lfloor\alpha^j\rfloor \geq
\frac{\alpha^j}2$ (yang sahih begitu $\alpha^j \geq 1$, yakni bagi setiap
$j \geq 0$: sebab $\lfloor x\rfloor \geq \frac x2$ bagi $x \geq
1$). Karena itu
\[
\sum_{j : k_j \geq n}\frac1{k_j^2}
\leq 4\sum_{j : \alpha^j \geq n}\alpha^{-2j}
\leq \frac{4}{1 - \alpha^{-2}}\cdot\frac1{n^2}
= \frac{C_\alpha}{n^2},
\]
(lewat deret geometrik dari $j$ pertama yang $\alpha^j \geq
n$). Lalu dengan menggabungkannya dengan pertanyaan 4--5, jumlah gandanya
berhingga; sehingga Borel--Cantelli (1), yang diterapkan bagi setiap
$\varepsilon$ rasional lalu diiriskan, memberikan
$\frac{S^*_{k_j} - \E S^*_{k_j}}{k_j} \to 0$ secara h.p., dan bersama
pertanyaan 3: $\frac{S^*_{k_j}}{k_j} \to m$ secara h.p.

\textbf{7.} Karena $Y_n \geq 0$, maka $n \mapsto S^*_n$
tak turun: sehingga bagi $k_j \leq n \leq k_{j+1}$,
\[
\frac{S^*_{k_j}}{k_{j+1}} \leq \frac{S^*_n}{n} \leq
\frac{S^*_{k_{j+1}}}{k_j},
\]
yang merupakan apitan yang ditampilkan itu setelah menyisipkan
$\frac{k_j}{k_{j+1}}$ dan $\frac{k_{j+1}}{k_j}$. Lalu karena
$\frac{k_{j+1}}{k_j} \to \alpha$, pertanyaan 6 memberikan secara h.p.
\[
\frac m\alpha \leq \liminf_n\frac{S^*_n}n \leq
\limsup_n\frac{S^*_n}n \leq \alpha m .
\]

\textbf{8.} Terapkanlah pertanyaan 7 bagi $\alpha = 1 + \frac1p$, dengan $p
\in \N^*$: yakni terbilang banyak kejadian h.p.; lalu pada
irisannya, dengan membiarkan $p \to \infty$: $\lim\frac{S^*_n}n =
m$ secara h.p. Bersama pertanyaan 1--3, $\frac{S_n}n \to \E X_1$
secara h.p.: yakni hukum kuat bilangan besar, di bawah \omterm{def:b3:probability:independence}{kebebasan}
berpasangan.

\textbf{9.} Hipotesis bertipe \omterm{def:b3:probability:independence}{kebebasan} muncul tiga
kali: (i) keaditifan variansnya (pertanyaan 5) ---
di sana berpasangan sudah cukup; (ii) \omterm{def:b3:probability:space}{distribusi} yang identik, pada
jumlah pemenggalannya (pertanyaan 2) dan perhitungan rata-ratanya
(pertanyaan 3) --- di sana tak ada \omterm{def:b3:probability:independence}{kebebasan} sama sekali; (iii)
Borel--Cantelli (1) (pertanyaan 2 dan 6) --- yang sahih tanpa
\omterm{def:b3:probability:independence}{kebebasan} apa pun. Jadi \omterm{def:b3:probability:independence}{kebebasan} bersama yang penuh tak pernah
dipanggil: itulah pengamatan Etemadi.

\textbf{10.} Tetapkanlah sebuah basis $b$ dan sebuah digit $r$. Maka digit
basis-$b$ $(d_k)$ milik $\omega$ yang seragam bersifat seragam i.i.d.\ pada
$\{0, \dots, b-1\}$ (sebab setiap nilai vektor digitnya menempati
selang berpanjang $b^{-m}$: yakni argumen
\cref{thm:b3:probability:existence} kata demi kata). Lalu hukum
kuatnya yang diterapkan pada variabel terbatas i.i.d.\ $\mathbf
1_{d_k = r}$ memberikan: secara h.p., frekuensi digit $r$
menuju $\frac1b$. Lalu dengan mengiriskan terbilang banyak
pasangan $(b, r)$: hampir setiap bilangan bersifat
\emph{normal sederhana pada setiap basis}. Adapun satu bilangan tak normal
yang eksplisit: $x = 0.100100100\ldots_2$ (dengan frekuensi angka satu
$\frac13 \neq \frac12$). Pertentangannya merendahkan hati: sebab hampir
semua bilangan normal, padahal bagi $\sqrt2$, $\eu$, atau $\pi$
kenormalannya tetap tak terbukti --- jadi teori \omterm{def:b3:measure:measure}{ukuran} mencacah
tanpa menunjukkan.

\textbf{11.} Menurut \cref{exo:b3:probability:10} yang diiterasi, satu
variabel seragam menghasilkan barisan \emph{vektor} seragam
i.i.d.\ $U_k$ pada $\intcc01^d$ (belahlah himpunan
digit setiap $U_n$ dari
\cref{thm:b3:probability:existence} menjadi $d$ subkeluarga).
Lalu bagi $g \in L^1(\intcc01^d)$, variabel $g(U_k)$ bersifat
i.i.d.\ dan \omterm{def:b3:lebesgue:l1}{terintegralkan} dengan rata-rata $\int g\,\dd\lambda_d$
(menurut transfernya): sehingga hukum kuatnya memberikan
\[
\frac1n\sum_{k=1}^ng(U_k)
\xrightarrow[n\to\infty]{\text{h.p.}}
\int_{\intcc01^d}g\,\dd\lambda_d :
\]
sehingga pengintegralan \omterm{ex:b3:probability:sllnapps}{Monte Carlo} konvergen hampir pasti, pada setiap
dimensi --- sedangkan besarnya galatnya adalah urusan teorema
limit pusatnya (\cref{ch:b3:clt}).

\textbf{12.} Ambillah $A_k = \{\abs{S_k} \geq \varepsilon\} \cap
\bigcap_{j<k}\{\abs{S_j} < \varepsilon\}$: maka $A_k$ saling
lepas dengan gabungan $A = \{\max_{k\leq n}\abs{S_k} \geq
\varepsilon\}$. Lalu
\[
\E S_n^2 \geq \sum_{k=1}^n\E\bigl[S_n^2\mathbf 1_{A_k}\bigr]
= \sum_{k=1}^n\E\Bigl[\bigl(S_k^2 + 2S_k(S_n - S_k) + (S_n
- S_k)^2\bigr)\mathbf 1_{A_k}\Bigr]
\geq \sum_{k=1}^n\E\bigl[S_k^2\mathbf 1_{A_k}\bigr],
\]
sebab suku silangnya lenyap: $S_k\mathbf 1_{A_k}$ merupakan
fungsi Borel koalisi $(Z_1, \dots, Z_k)$, yang
\omterm{def:b3:probability:independence}{bebas} dari $S_n - S_k$, yakni fungsi $(Z_{k+1},
\dots, Z_n)$ (\cref{thm:b3:probability:independence}), sehingga
$\E[S_k\mathbf 1_{A_k}(S_n - S_k)] = \E[S_k\mathbf
1_{A_k}]\,\E[S_n - S_k] = 0$. Sedangkan pada $A_k$, $S_k^2 \geq
\varepsilon^2$, sehingga $\E S_n^2 \geq
\varepsilon^2\sum_k\P(A_k) = \varepsilon^2\P(A)$; dan $\E
S_n^2 = \sum_{k\leq n}\V(Z_k)$ (sebab variansnya menjumlah). Adapun
langkah yang menentukan adalah pemfaktorannya: sebab $S_k\mathbf 1_{A_k}$ merupakan
fungsi \emph{taklinear} seluruh blok pertamanya, dan
\omterm{def:b3:probability:independence}{kebebasannya} dari blok kedua adalah \omterm{def:b3:probability:independence}{kebebasan} koalisi ---
sedangkan \omterm{def:b3:probability:independence}{kebebasan} berpasangan $Z_i$ hanya
menghilangkan korelasi pasangan dan takkan membenarkannya.

\textbf{13.} Tetapkanlah $N$ lalu terapkan pertanyaan 12 pada $Z_{N+1},
\dots, Z_{N+m}$:
\[
\P\Bigl(\max_{N < k \leq N+m}\abs{S_k - S_N} >
\varepsilon\Bigr) \leq
\frac1{\varepsilon^2}\sum_{j=N+1}^{N+m}\V(Z_j) \leq
\frac{r_N}{\varepsilon^2},
\qquad r_N = \sum_{j>N}\V(Z_j) .
\]
Kejadiannya naik terhadap $m$; sehingga \omterm{def:b3:topology:continuity}{kekontinuan} dari bawah memberikan
$\P(\sup_{k>N}\abs{S_k - S_N} > \varepsilon) \leq
r_N/\varepsilon^2$, sedangkan $r_N \to 0$ menurut hipotesisnya. Karena itu
bagi setiap $p \in \N^*$, $\P\bigl(\bigcap_N\{\sup_{k>N}
\abs{S_k - S_N} > \frac1p\}\bigr) \leq \inf_Np^2r_N = 0$:
jadi hampir pasti, bagi setiap $p$ ada $N$ dengan
$\sup_{k>N}\abs{S_k - S_N} \leq \frac1p$ (iriskanlah
terbilang banyak kejadian h.p.\ atas $p$), sehingga $\abs{S_k -
S_l} \leq \frac2p$ bagi setiap $k, l > N$: jadi jumlah parsialnya
Cauchy secara h.p., sehingga konvergen secara h.p.

\textbf{14.} Variabel $Z_n = x_n\varepsilon_n$ bersifat
\omterm{def:b3:probability:independence}{bebas} (sebab fungsi Borel variabel yang \omterm{def:b3:probability:independence}{bebas},
\cref{exo:b3:probability:3}(a)), terpusat, dengan $\V(Z_n) =
x_n^2$: sehingga pertanyaan 13 berlaku bila $\sum_nx_n^2 < \infty$ dan
memberikan konvergensi h.p. Secara umum, bagi setiap $N$
konvergensi $\sum_nx_n\varepsilon_n$ tak terpengaruh oleh
nilai $\varepsilon_1, \dots, \varepsilon_N$: sehingga
kejadian konvergensinya terletak pada $\sigma$-aljabar ekor
barisan \omterm{def:b3:probability:independence}{bebas} $(\varepsilon_n)$, jadi \omterm{thm:b3:probability:zeroone}{hukum nol--satu} Kolmogorov
(\cref{thm:b3:probability:zeroone}) memaksa
peluangnya menjadi $0$ atau $1$.

\textbf{15.} (a) Dengan membelahnya pada aras $\theta\E Z$ lalu
memakai Cauchy--Schwarz pada keping atasnya,
\[
\E Z = \E\bigl[Z\mathbf 1_{Z \leq \theta\E Z}\bigr] +
\E\bigl[Z\mathbf 1_{Z > \theta\E Z}\bigr]
\leq \theta\,\E Z + \sqrt{\E Z^2}\,
\sqrt{\P(Z > \theta\E Z)} ,
\]
sehingga $(1 - \theta)\E Z \leq \sqrt{\E Z^2\,\P(Z > \theta\E
Z)}$; lalu kuadratkanlah. (b) Uraikanlah $T_n^4 =
\sum_{i,j,k,l}x_ix_jx_kx_l\,
\E[\varepsilon_i\varepsilon_j\varepsilon_k\varepsilon_l]$:
\omterm{def:b3:probability:space}{nilai harapannya} $1$ bila indeksnya berpasangan (yakni keempatnya
sama, atau dua pasangan berbeda, yang terakhir dalam $3$
susunan) dan $0$ selainnya (sebab tanda yang tak berpasangan berrata-rata
nol dan terfaktorkan keluar oleh \omterm{def:b3:probability:independence}{kebebasannya}). Karena itu
\[
\E T_n^4 = \sum_kx_k^4 + 3\sum_{i\neq j}x_i^2x_j^2 =
3s_n^4 - 2\sum_kx_k^4 \leq 3s_n^4 .
\]
(c) Paley--Zygmund dengan $Z = T_n^2$, $\E Z = s_n^2$,
$\theta = \frac14$:
\[
\P\Bigl(\abs{T_n} > \frac{s_n}2\Bigr) = \P\Bigl(T_n^2 >
\frac{s_n^2}4\Bigr) \geq \Bigl(\frac34\Bigr)^2
\frac{s_n^4}{3s_n^4} = \frac3{16} .
\]
Seandainya deretnya konvergen dengan peluang positif, maka ia
akan konvergen secara h.p.\ (pertanyaan 14), sehingga $\sup_n\abs{T_n} <
\infty$ secara h.p., dan suatu $M$ akan memenuhi
$\P(\sup_n\abs{T_n} > M) < \frac3{16}$; padahal begitu $s_n
> 2M$, $\P(\abs{T_n} > M) \geq \P(\abs{T_n} > \frac{s_n}2)
\geq \frac3{16}$: yakni kontradiksi. Jadi divergensinya hampir
pasti, dan bersama pertanyaan 14 dikotominya \omterm{def:b3:complete:complete}{lengkap}.

\textbf{16.} Di sini $x_n = n^{-s}$ dan $\sum_nn^{-2s} <
\infty$ tepat saat $s > \frac12$: jadi menurut pertanyaan 14--15,
$\sum_n\frac{\varepsilon_n}{n^s}$ konvergen secara h.p.\ jika dan
hanya jika $s > \frac12$ (sedangkan bagi $s \leq \frac12$, ia divergen
secara h.p.). Untuk $\frac12 < s \leq 1$ konvergensinya
tak pernah mutlak. Perbandingannya mendidik: sebab tanda yang berselang-seling
\omterm{prop:b3:galois:perfect}{sempurna} meniadakan diri sekuat $n^{-s}$ bagi setiap $s
> 0$, sedangkan tanda acak yang khas meniadakan diri hanya sekuat
akar kuadrat --- sebab jalan acak pertanyaan 21
tumbuh seperti $\sqrt n$, dan penjumlahan Abel mengubah persis
pertumbuhan itu menjadi konvergensi $\sum\varepsilon_nn^{-s}$
bagi $s > \frac12$.

\textbf{17.} (a) Berlaku $\cosh\lambda =
\sum_k\frac{\lambda^{2k}}{(2k)!}$ dan $\eu^{\lambda^2/2} =
\sum_k\frac{\lambda^{2k}}{2^kk!}$; sedangkan $(2k)! \geq 2^kk!$
berlaku suku demi suku, sebab $\frac{(2k)!}{k!} =
\prod_{i=1}^k(k + i) \geq \prod_{i=1}^k(2i) = 2^kk!$ (karena setiap
faktornya memenuhi $k + i \geq 2i$ bagi $i \leq k$), sehingga
sesungguhnya $(2k)! \geq 2^k(k!)^2 \geq 2^kk!$. (b) Perhatikanlah $a \leq 0 \leq b$ (sebab $Z$
terpusat), lalu menurut kecembungan $z \mapsto \eu^{\lambda z}$,
bagi $z \in \intcc ab$:
\[
\eu^{\lambda z} \leq \frac{b - z}{b - a}\,\eu^{\lambda a} +
\frac{z - a}{b - a}\,\eu^{\lambda b},
\qquad\text{sehingga}\qquad
\E\,\eu^{\lambda Z} \leq \frac{b\,\eu^{\lambda a} -
a\,\eu^{\lambda b}}{b - a}
= (1 - p)\eu^{-pt} + p\,\eu^{(1-p)t} = \eu^{\varphi(t)}
\]
dengan $p = \frac{-a}{b-a} \in \intcc01$, $t = \lambda(b -
a)$, $\varphi(t) = -pt + \log(1 - p + p\eu^t)$. Maka
$\varphi(0) = 0$, $\varphi'(t) = -p + \frac{p\eu^t}{1 - p +
p\eu^t}$ lenyap di $0$, dan $\varphi''(t) = \rho(1 -
\rho) \leq \frac14$ bagi $\rho = \frac{p\eu^t}{1 - p +
p\eu^t} \in \intcc01$: sehingga Taylor pada orde $2$ memberikan
$\varphi(t) \leq \frac{t^2}8 = \frac{\lambda^2(b-a)^2}8$.

\textbf{18.} Untuk $\lambda > 0$, Markov yang diterapkan pada
variabel positif $\eu^{\lambda(S_n - \E S_n)}$
(\cref{prop:b3:probability:markov}) dan rumus hasil kali
bagi variabel yang \omterm{def:b3:probability:independence}{bebas} memberikan
\[
\P(S_n - \E S_n \geq t) \leq \eu^{-\lambda
t}\prod_{i=1}^n\E\,\eu^{\lambda(X_i - \E X_i)}
\leq \exp\Bigl(-\lambda t +
\frac{\lambda^2}8\sum_i(b_i - a_i)^2\Bigr),
\]
menurut pertanyaan 17(b) yang diterapkan pada setiap $X_i - \E X_i
\in \intcc{a_i - \E X_i}{b_i - \E X_i}$ yang terpusat (dengan lebar yang sama).
Lalu meminimumkan eksponennya di $\lambda = \frac{4t}{D}$, dengan $D =
\sum_i(b_i - a_i)^2$, menghasilkan $-\frac{2t^2}D$. Sedangkan ekor
bawahnya menyusul dengan menerapkan hasilnya pada $(-X_i)$.

\textbf{19.} Ambillah $t = n\varepsilon$ dan $D = n(b - a)^2$:
\[
\P\Bigl(\Bigl|\frac{S_n}n - m\Bigr| \geq \varepsilon\Bigr)
\leq 2\exp\Bigl(\frac{-2n^2\varepsilon^2}{n(b-a)^2}\Bigr)
= 2\exp\Bigl(\frac{-2n\varepsilon^2}{(b-a)^2}\Bigr),
\]
yang terjumlahkan terhadap $n$ (yakni deret bertipe geometrik):
sehingga Borel--Cantelli (\cref{thm:b3:probability:borelcantelli})
memberikan bahwa secara h.p.\ $\abs{\frac{S_n}n - m} < \varepsilon$
pada akhirnya; lalu mengiriskannya atas $\varepsilon = \frac1p$
menghasilkan $\frac{S_n}n \to m$ secara h.p. Perbandingannya: Etemadi meminta
hanya $X_1 \in L^1$ dan \omterm{def:b3:probability:independence}{kebebasan} berpasangan, lalu tak memberikan
laju; sedangkan Hoeffding meminta keterbatasan dan \omterm{def:b3:probability:independence}{kebebasan} penuh,
lalu memberikan jaminan eksponensial yang eksplisit pada setiap
$n$ berhingga --- jadi kedua teoremanya menjawab pertanyaan yang berbeda
tentang limit yang sama.

\textbf{20.} Variabel $g(U_k)$ bersifat i.i.d.\ dengan nilai di
$\intcc01$ dan berrata-rata $\int g\,\dd\lambda_d$ (menurut transfernya), sehingga
pertanyaan 18 dengan $b_i - a_i = 1$, $t = n\varepsilon$ memberikan
batas dua sisinya $2\eu^{-2n\varepsilon^2} \leq \delta$
begitu $\eu^{2n\varepsilon^2} \geq \frac2\delta$, yakni
$n \geq \frac{\log(2/\delta)}{2\varepsilon^2}$. Untuk
$\varepsilon = \delta = 10^{-2}$:
\[
n \geq \frac{\log 200}{2\cdot10^{-4}} =
\frac{5.2983\ldots}{0.0002} \approx 26\,492 :
\]
jadi sekitar $26\,500$ cuplikan menjamin ketelitian $1\%$ dengan
keyakinan $99\%$ --- pada setiap dimensi $d$, bagi setiap
integran \omterm{def:b3:lebesgue:measurable}{terukur} yang bernilai di $\intcc01$. Sedangkan hukum kuat
pertanyaan 11 menjanjikan konvergensi tanpa jaminan pada $n$ berhingga;
lalu kisi deterministik dengan $k$ titik per sumbu
memakan $k^d$ evaluasi, yang eksponensial terhadap $d$. Jadi pemusatannyalah
yang menjadikan \omterm{ex:b3:probability:sllnapps}{Monte Carlo} sebuah \emph{metode} alih-alih sekadar
angan-angan.

\textbf{21.} Lewat \omterm{def:b3:probability:independence}{kebebasan} dan rumus hasil kalinya:
$\E\,\eu^{\lambda S_n} = (\E\,\eu^{\lambda\varepsilon_1})^n
= (\cosh\lambda)^n \leq \eu^{n\lambda^2/2}$ menurut pertanyaan
17(a). Lalu Markov pada $\eu^{\lambda S_n}$:
\[
\P(S_n \geq x) \leq \eu^{-\lambda x + n\lambda^2/2}
= \eu^{-x^2/(2n)}
\qquad\text{pada optimumnya } \lambda = \frac xn,
\]
sedangkan batas setangkupnya bagi $-S_n$ (yang berdistribusi sama) menggandakan
konstantanya bagi $\abs{S_n}$.

\textbf{22.} Tetapkanlah $\eta > 0$ lalu tetapkan $x_n = (1 +
\eta)\sqrt{2n\log n}$ bagi $n \geq 2$:
\[
\P(\abs{S_n} \geq x_n) \leq 2\exp\bigl(-(1 +
\eta)^2\log n\bigr) = \frac{2}{n^{(1+\eta)^2}},
\]
yang terjumlahkan sebab $(1 + \eta)^2 > 1$. Lalu Borel--Cantelli: secara h.p.\
$\abs{S_n} < (1 + \eta)\sqrt{2n\log n}$ bagi setiap $n$ yang besar,
sehingga $\limsup_n\frac{\abs{S_n}}{\sqrt{2n\log n}} \leq 1 +
\eta$ secara h.p.; lalu mengiriskan kejadian h.p.-nya bagi $\eta =
\frac1p$, $p \in \N^*$, memberikan klaimnya. Jadi jalan acak berukuran
$n$ beramplitudo khas $\sqrt n$ (yakni variansnya), sedangkan
bahkan gerak jauhnya yang terburuk melampaui skala itu paling banyak sebesar
$\sqrt{2\log n}$.

\textbf{23.} Dengan $n_j = 2^j$ dan $x = (1 +
\eta)\sqrt{2n_j\log\log n_j}$ (yang terdefinisi bagi $j \geq 2$),
pertanyaan 21 memberikan
\[
\P\bigl(S_{n_j} \geq x\bigr) \leq \exp\bigl(-(1 +
\eta)^2\log\log n_j\bigr) = (j\log 2)^{-(1+\eta)^2},
\]
yang terjumlahkan terhadap $j$ sebab $(1 + \eta)^2 > 1$: sehingga Borel--Cantelli
dan $\eta = \frac1p$ memberikan $\limsup_jS_{n_j}/\sqrt{2n_j
\log\log n_j} \leq 1$ secara h.p. Adapun yang kurang bagi paruh atasnya
yang penuh adalah jembatan antara pos pemeriksaannya: sebab orang harus
menunjukkan bahwa $\max_{n_j\leq n\leq n_{j+1}}S_n$ melampaui
$(1+\eta)\sqrt{2n_j\log\log n_j}$ hanya berhingga kali,
dan itu menuntut ketaksamaan maksimal berekor \emph{Gauss}
(yakni ketaksamaan pantulan Lévy atau ketaksamaan Ottaviani,
yang tak dibuktikan di sini). Sedangkan pertanyaan 12 secara kuantitatif
terlalu lemah: sebab ia membatasi peluangnya oleh
\[
\frac{n_j}{(1+\eta)^2\,2n_j\log\log n_j}
= \frac{1}{2(1+\eta)^2\log(j\log2)},
\]
yang menuju $0$ tetapi \emph{tak terjumlahkan} terhadap $j$:
sehingga Borel--Cantelli tak dapat menyimpulkan. Adapun paruh bawah hukum
logaritma teriterasinya menerapkan lema Borel--Cantelli
kedua pada pertambahan yang \omterm{def:b3:probability:independence}{bebas}
$S_{n_{j+1}} - S_{n_j}$, dengan memakai batas bawah yang sepadan bagi
ekor bertipe Gauss. Kedua pertajaman itu merupakan peluang Tahun 3
yang sejati, satu kuliah lebih jauh; sedangkan yang diberikan soal ini
tanpa bantuan adalah skala logaritma-teriterasi yang tepat
sepanjang waktu geometriknya.

\textbf{24.} Setiap $\hat p_i$ merupakan rata-rata $n$ variabel
indikator i.i.d.\ yang bernilai di $\intcc01$ dan berrata-rata
$\P(A_i)$: sehingga Hoeffding memberikan $\P(\abs{\hat p_i - \P(A_i)} >
\varepsilon) \leq 2\eu^{-2n\varepsilon^2}$. Lalu batas gabungannya
mengalikannya dengan $N$. Lalu menyelesaikan $2N\eu^{-2n\varepsilon^2} \leq
\delta$: $n \geq \frac{\ln(2N/\delta)}{2\varepsilon^2}$.
Secara numerik: $\ln\frac{2\cdot10^6}{0.05} =
\ln(4\cdot10^7) \approx 17.5$, sehingga $n \geq
\frac{17.5}{2\cdot10^{-4}} \approx 87\,600$: jadi menaksir
\emph{satu} peluang sampai $\pm1\%$ memakan sekitar $18\,500$
cuplikan (yakni $\ln(2/\delta)/2\varepsilon^2$), sedangkan
\emph{sejuta} peluang hanya $\approx 4.7$ kali lebih banyak
--- jadi keseragamannya berharga $\ln N$, bukan $N$: yakni pengamatan yang
menjadikan peminimuman risiko empiris, dan bersamanya pembelajaran
mesin, mungkin secara statistik.

\textbf{25.} Variabel $X_n = \frac{\varepsilon_n}
{n^\alpha}$ bersifat \omterm{def:b3:probability:independence}{bebas}, terpusat, terbatas, dengan
$\sum_n\V(X_n) = \sum_nn^{-2\alpha}$. Jika $\alpha >
\frac12$: maka deret variansnya konvergen, sehingga
teorema satu deretnya (Bagian VI) memberikan konvergensi h.p.\
$\sum X_n$. Sedangkan jika $\alpha \leq \frac12$: deret variansnya
divergen, sehingga paruh konversnya (yakni argumen
Paley--Zygmund Bagian VI, yang berlaku sebab sukunya
terbatas oleh $1$) memberikan divergensi h.p. Adapun konvergensi
mutlaknya meminta $\sum n^{-\alpha} < \infty$: yakni $\alpha > 1$.
Jadi pada $\intoc{\frac12}1$, deretnya konvergen secara h.p.\ walaupun
$\sum\abs{X_n} = \infty$ secara pasti: sebab tandanya bersekongkol untuk
meniadakan diri, dengan peluang satu --- yakni konvergensi lewat
peniadaan, yang tak terlihat oleh uji mutlak mana pun, dan (menurut
\omterm{thm:b3:probability:zeroone}{hukum nol--satunya}) dengan vonis yang deterministik pula.

\end{solution}
